% formal/product-triangle-bounce-classification.tex
%
% Developer-facing proof packet for the two- versus three-bounce comparison
% for a Lagrangian product of two triangles.  The new mathematics is
% agent-derived and awaits Jörn's review.

\subsection{Status, scope, and conventions}
\label{subsec:product-triangle-bounce-status}

This packet compares the complete mathematical two-bounce value with the
transition-admissible all-singleton three-block words for a product of two
triangles.  It also proves, by an explicit placement construction, that every
strict directed Hamiltonian word in this setting is realized by a proper
strong three-bounce billiard.  It does not infer global placement for general
words merely from local transition tests.

\begin{unverified}
Every new lemma, proposition, theorem, and corollary in this packet is an
agent-derived proof at agent-review status.  None has yet been accepted by
Jörn.  The exact computation cited at the end is supporting falsifier evidence,
not part of any proof.
\end{unverified}

Let
\[
 P=\{q\in\mathbb R^2:\langle a_i,q\rangle\leq1,
                         \ i\in\mathbb Z/3\},\qquad
 Q=\{p\in\mathbb R^2:\langle b_j,p\rangle\leq1,
                         \ j\in\mathbb Z/3\}
\]
be full-dimensional triangles containing the origin in their interiors, with
irredundant displayed inequalities.  Their unique normalized positive
dependencies are
\begin{equation}
\label{eq:product-triangle-positive-dependencies}
 \alpha_i>0,\quad \sum_i\alpha_i=1,\quad
 \sum_i\alpha_i a_i=0,
 \qquad
 \gamma_j>0,\quad \sum_j\gamma_j=1,\quad
 \sum_j\gamma_j b_j=0.
\end{equation}
Put
\[
 A_i=\alpha_i a_i,\qquad B_j=\gamma_j b_j,
 \qquad c_{ij}=\langle A_i,B_j\rangle,
 \qquad M=\max_{i,j}|c_{ij}|.
\]
Thus every row and every column of \(C=(c_{ij})\) sums to zero.

The product coordinates and symplectic convention are
\[
 (q_1,q_2,p_1,p_2),\qquad
 \omega_0((a,0),(0,b))=\langle a,b\rangle,
 \qquad
 \omega_0((0,b),(a,0))=-\langle a,b\rangle.
\]
For an active Reeb-traversal word
\(
 \sigma=(a_0,b_0,a_1,b_1,a_2,b_2)
\), the project objective is
\begin{equation}
\label{eq:product-triangle-active-word-objective}
 Q_\sigma(\beta)=
 \sum_{1\leq j<k\leq6}\beta_j\beta_k
 \omega_0(\widehat a_{\sigma_j},\widehat a_{\sigma_k}),
\end{equation}
where \(\widehat a_i=(a_i,0)\) and
\(\widehat b_j=(0,b_j)\).  The action convention is
\begin{equation}
\label{eq:product-triangle-action-convention}
 \mathcal A_\sigma=\frac1{2\max Q_\sigma}.
\end{equation}
This is the active-word convention audited in
Subsection~\ref{subsec:hk2017-active-word-qp-convention}; the fixed word in
Haim--Kislev's displayed orientation is its reverse.

The complete mathematical two-bounce action is the value
\(A_2(P,Q)\) of Definition~\ref{def:product-two-bounce-a2}, equivalently
\(1/(2Q_{2,\max})\) by
Theorem~\ref{thm:product-two-bounce-complete-qp}.  We do not use the
algorithmic pruning labels \(A2\) and \(A3\) for these actions.

\subsection{The triangle difference-body polar}
\label{subsec:product-triangle-difference-polar}

\begin{unverified}
\begin{lemma}[Difference-body polar of a triangle]
\label{lem:product-triangle-difference-polar}
With the notation above,
\[
 (P-P)^\circ=\operatorname{conv}\{\pm A_0,\pm A_1,\pm A_2\},
 \qquad
 (Q-Q)^\circ=\operatorname{conv}\{\pm B_0,\pm B_1,\pm B_2\}.
\]
Consequently,
\begin{equation}
\label{eq:product-triangle-a2-m}
 Q_{2,\max}(P,Q)=\frac M2,
 \qquad A_2(P,Q)=\frac1M.
\end{equation}
\end{lemma}

\begin{proof}
Define the slack coordinates
\[
 \lambda_i(q)=\alpha_i(1-\langle a_i,q\rangle).
\]
Equation~\eqref{eq:product-triangle-positive-dependencies} gives
\(\sum_i\lambda_i(q)=1\).  The map
\(q\mapsto(\lambda_0(q),\lambda_1(q),\lambda_2(q))\) has rank two and
is therefore an affine bijection from \(\mathbb R^2\) to the affine plane
\(\sum_i\lambda_i=1\).  Nonnegativity of its coordinates is exactly the
three inequalities defining \(P\), so it restricts to an affine bijection
from \(P\) to the standard two-simplex \(\Delta_2\).

The difference set of that simplex is
\[
 \Delta_2-\Delta_2
 =\{\delta\in\mathbb R^3:\sum_i\delta_i=0,
                              \ |\delta_i|\leq1\text{ for every }i\}.
\]
Indeed, the stated conditions imply that the sum of the positive coordinates
is at most one: if two are positive, their sum is the absolute value of the
sole negative coordinate; if one is positive, the assertion is immediate.
The positive and negative parts of \(\delta\), supplemented by the same
unused simplex mass, then exhibit \(\delta\) as a difference of two points of
\(\Delta_2\).  The reverse inclusion is immediate.

For \(d=q-q'\),
\[
 \lambda_i(q)-\lambda_i(q')=-\langle A_i,d\rangle.
\]
Since \(\sum_i A_i=0\), the preceding characterization yields
\[
 d\in P-P
 \quad\Longleftrightarrow\quad
 |\langle A_i,d\rangle|\leq1\quad(i=0,1,2).
\]
The right-hand side is the polar of
\(\operatorname{conv}\{\pm A_i\}\).  Taking polars proves the identity for
\(P\); the proof for \(Q\) is identical.

A bilinear functional on the product of two polytopes attains its maximum on
generating vertices.  Hence
\[
 \max_{x\in(P-P)^\circ,\ y\in(Q-Q)^\circ}\langle x,y\rangle
 =\max_{i,j}|\langle A_i,B_j\rangle|=M.
\]
The complete two-bounce formula
\eqref{eq:product-two-bounce-complete-qp} now gives
\eqref{eq:product-triangle-a2-m}.
\end{proof}
\end{unverified}

\subsection{One directed Hamiltonian word}
\label{subsec:product-triangle-one-word}

The directed transition conditions for the explicitly oriented word
\[
 \sigma=(a_0,b_0,a_1,b_1,a_2,b_2)
\]
are
\begin{equation}
\label{eq:product-triangle-directed-signs}
 \langle a_r,b_r\rangle\geq0,
 \qquad
 \langle a_{r+1},b_r\rangle\leq0
 \qquad(r\in\mathbb Z/3).
\end{equation}
Indeed, these are respectively
\(\omega_0((a_r,0),(0,b_r))\geq0\) and
\(\omega_0((0,b_r),(a_{r+1},0))\geq0\).
Cross-factor product facets always meet, so
Corollary~\ref{cor:ridge-sufficiency} leaves precisely these signs as the
local transition conditions.  Call the word \emph{strict} when all six
inequalities are strict.

Set
\begin{equation}
\label{eq:product-triangle-pn}
 p_r=c_{rr},\qquad n_r=-c_{r+1,r}.
\end{equation}
Under \eqref{eq:product-triangle-directed-signs}, all \(p_r,n_r\) are
nonnegative; under strict signs they are positive.

\begin{unverified}
\begin{lemma}[Hamiltonian sign-cell algebra]
\label{lem:product-triangle-sign-cell-algebra}
Under the weak signs \eqref{eq:product-triangle-directed-signs}, there is a
common scalar \(S>0\) such that
\begin{equation}
\label{eq:product-triangle-common-s}
 S=p_0+n_1=p_1+n_2=p_2+n_0.
\end{equation}
Every positive feasible weight vector for \(\sigma\) is uniquely of the form
\begin{equation}
\label{eq:product-triangle-word-weights}
 \beta_{a_r}=s\alpha_r,\qquad
 \beta_{b_r}=(1-s)\gamma_r,\qquad 0<s<1,
\end{equation}
and its objective is
\begin{equation}
\label{eq:product-triangle-q-s}
 Q_\sigma(s)=2s(1-s)S.
\end{equation}
Thus
\begin{equation}
\label{eq:product-triangle-a3-s}
 \max Q_\sigma=\frac S2,\qquad
 A_{3,\sigma}(P,Q):=\frac1{2\max Q_\sigma}=\frac1S.
\end{equation}
\end{lemma}

\begin{proof}
The zero sum of column \(r\) gives its unused entry as
\[
 c_{r-1,r}=n_r-p_r.
\]
Substitution in the three row-sum equations gives, cyclically,
\[
 p_0+n_1=p_1+n_2,\qquad
 p_1+n_2=p_2+n_0,\qquad
 p_2+n_0=p_0+n_1,
\]
which proves \eqref{eq:product-triangle-common-s}.  If \(S=0\),
nonnegativity would make all six \(p_r,n_r\) zero, and then the column formula
would make \(C=0\).  Since the \(A_i\) span \(\mathbb R^2\), this would force
every \(B_j=0\), contrary to the facet normalization.  Hence \(S>0\).

Factorwise closure and uniqueness of each positive dependence in
\eqref{eq:product-triangle-positive-dependencies} force the \(a\)-weights to
be a positive multiple of \((\alpha_r)_r\), and independently the
\(b\)-weights to be a positive multiple of \((\gamma_r)_r\).  Normalization
of the total weight gives exactly \eqref{eq:product-triangle-word-weights}.

In the active order \(A_0,B_0,A_1,B_1,A_2,B_2\), the pair \(A_r,B_j\)
contributes with sign \(+\) when \(r\leq j\), and with sign \(-\) when
\(j<r\).  Therefore
\[
 \begin{aligned}
 Q_\sigma(s)
 &=s(1-s)\left(\sum_{r\leq j}c_{rj}
                         -\sum_{j<r}c_{rj}\right)\\
 &=2s(1-s)\sum_{r\leq j}c_{rj}.
 \end{aligned}
\]
The second equality uses \(\sum_{r,j}c_{rj}=0\).  In the upper-triangular
sum, row zero contributes its full row sum, row one contributes
\(c_{11}+c_{12}=-c_{10}=n_0\), and row two contributes \(c_{22}=p_2\).
Thus the sum is \(n_0+p_2=S\), proving
\eqref{eq:product-triangle-q-s}.  Finally,
\(s(1-s)\leq1/4\), with equality uniquely at \(s=1/2\).  Combining this
with the action convention \eqref{eq:product-triangle-action-convention}
proves \eqref{eq:product-triangle-a3-s}.
\end{proof}
\end{unverified}

\subsection{Strict comparison and the complete weak boundary}
\label{subsec:product-triangle-comparison}

\begin{unverified}
\begin{theorem}[Wordwise two- versus three-bounce comparison]
\label{thm:product-triangle-wordwise-comparison}
Let \(\sigma\) satisfy the six weak directed signs
\eqref{eq:product-triangle-directed-signs}.  Then
\[
 A_{3,\sigma}(P,Q)\leq A_2(P,Q).
\]
More precisely,
\begin{enumerate}
  \item if all six traversed pairings are nonzero, then
  \(A_{3,\sigma}(P,Q)<A_2(P,Q)\);
  \item equality holds if and only if at least one of the six traversed
  pairings is zero.
\end{enumerate}
A zero in one of the three unused chord pairings \(c_{r-1,r}\) does not
remove strictness when all six traversed pairings are nonzero.
\end{theorem}

\begin{proof}
Equation~\eqref{eq:product-triangle-common-s} can be written cyclically as
\[
 S=p_r+n_{r+1}=p_{r-1}+n_r.
\]
Thus \(0\leq p_r,n_r\leq S\), and the unused column entry obeys
\[
 |c_{r-1,r}|=|n_r-p_r|\leq S.
\]
It follows that \(M\leq S\).  Lemma~\ref{lem:product-triangle-difference-polar}
and Lemma~\ref{lem:product-triangle-sign-cell-algebra} give
\[
 A_{3,\sigma}=\frac1S\leq\frac1M=A_2.
\]

If all six traversed pairings are nonzero, then every \(p_r,n_r\) is
positive.  The two expressions for \(S\) show \(p_r,n_r<S\), and hence
\(|n_r-p_r|<S\).  Therefore every entry of \(C\) has absolute value
strictly below \(S\), so \(M<S\) and the action inequality is strict.

Conversely, if \(p_r=0\), then \(n_{r+1}=S\); if \(n_r=0\), then
\(p_{r-1}=S\).  In either case \(M=S\) and the actions are equal.  The
previous paragraph proves the converse implication.  In particular, an
unused chord equality \(c_{r-1,r}=n_r-p_r=0\) has no effect on strictness
when all \(p_r,n_r\) remain positive.
\end{proof}
\end{unverified}

\subsection{Finite triangle classification}
\label{subsec:product-triangle-finite-classification}

For each factor, a \(k=3\) choice in the complete alternating block family
consists of three pairwise disjoint nonempty blocks drawn from only three
facets.  Therefore every block is a singleton and all three facets occur.
Consequently every \(k=3\) word is an alternating Hamiltonian six-cycle; no
two-facet block can occur in this class.

Let \(\mathcal H(P,Q)\) be the finite set of these oriented Hamiltonian words
which satisfy all six weak directed transition signs.  Relabelling each word
cyclically puts it in the form used above.  Define unambiguously
\[
 A_3^{\mathrm{dir}}(P,Q)
 =\min_{\sigma\in\mathcal H(P,Q)}A_{3,\sigma}(P,Q),
\]
with value \(+\infty\) if \(\mathcal H(P,Q)=\varnothing\).  This is the
complete transition-admissible all-singleton three-block value; it is not a
claim that arbitrary local transition cycles for arbitrary products have a
global physical placement.

\begin{unverified}
\begin{corollary}[Finite two- versus three-block classification]
\label{cor:product-triangle-finite-classification}
For two triangles,
\begin{equation}
\label{eq:product-triangle-global-finite-classification}
 \min\{A_2(P,Q),A_3^{\mathrm{dir}}(P,Q)\}<A_2(P,Q)
\end{equation}
if and only if the cross-pairing graph has a strict directed Hamiltonian
six-cycle.  If no strict cycle exists, the left-hand minimum equals \(A_2\):
every weakly directed three-block word either is absent or has an incident
zero and therefore has action exactly \(A_2\).

On the generic stratum where all nine cross pairings are nonzero, every weak
directed Hamiltonian cycle is strict.  Hence on that stratum absence of a
strict cycle is equivalent to absence of any transition-admissible
three-block word.
\end{corollary}

\begin{proof}
The block-count argument above exhausts the \(k=3\) choices.  Apply
Theorem~\ref{thm:product-triangle-wordwise-comparison} to each member of the
finite set \(\mathcal H(P,Q)\).  A member has action below \(A_2\) exactly
when all six of its traversed pairings are nonzero, which under its weak signs
is exactly strictness.  If no member is strict, every member has an incident
zero and action \(A_2\); the explicit \(A_2\) term still supplies the
two-bounce value when the set is empty.  The final assertion is immediate
when none of the nine pairings vanishes.
\end{proof}
\end{unverified}

This corollary is a complete classification of the stated finite QP
comparison, including its zero boundary.  By itself it does not classify
degenerate strong billiards with vertex contacts, nor does it prove that local
transition feasibility gives global physical placement.  The next proposition
supplies the latter implication precisely for the strict triangle cycles used
in the strict half of the corollary.

\subsection{Strict cycles are genuine strong three-bounce billiards}
\label{subsec:product-triangle-physical-realization}

Write \(F_P(a_i)=\{q\in P:\langle a_i,q\rangle=1\}\), and similarly for
\(F_Q(b_j)\).  We first isolate the placement step used in both factors.

\begin{unverified}
\begin{lemma}[Homothetic side placement]
\label{lem:product-triangle-homothetic-placement}
Let \(H=\{x:\langle u_i,x\rangle\leq1,\ i\in\mathbb Z/3\}\) be a
full-dimensional triangle, with normalized positive dependence
\(\sum_i\lambda_i u_i=0\), \(\sum_i\lambda_i=1\).  Let \(x_0,x_1,x_2\)
be the vertices of a nondegenerate triangle \(T\), and suppose \(u_i\) is
strictly maximized on \(T\) at \(x_i\).  Then there are unique \(z\in\mathbb
R^2\) and \(t>0\) such that
\[
 z+t x_i\in\operatorname{relint}F_H(u_i)
 \qquad(i\in\mathbb Z/3).
\]
Moreover,
\begin{equation}
\label{eq:product-triangle-placement-scale}
 t=\left(\sum_i\lambda_i\langle u_i,x_i\rangle\right)^{-1}.
\end{equation}
\end{lemma}

\begin{proof}
The denominator in \eqref{eq:product-triangle-placement-scale} is invariant
under translating all \(x_i\), because \(\sum_i\lambda_i u_i=0\).  Translate
\(T\) so that the origin is in its interior.  Every strict supporting value
\(\langle u_i,x_i\rangle=h_T(u_i)\) is then positive, so the denominator is
positive.

We seek \(z,t\) satisfying
\(\langle u_i,z+t x_i\rangle=1\) for all \(i\).  The image of
\(z\mapsto(\langle u_i,z\rangle)_i\) is exactly the two-plane
\(\sum_i\lambda_i y_i=0\).  Thus the three equations are compatible exactly
when
\[
 0=\sum_i\lambda_i(1-t\langle u_i,x_i\rangle),
\]
which gives the displayed positive and unique \(t\); the spanning property of
the \(u_i\) then gives a unique \(z\).  Finally, if \(k\neq i\), strict support
gives \(\langle u_k,x_i\rangle<\langle u_k,x_k\rangle\).  Comparing the
two placed points shows
\(\langle u_k,z+t x_i\rangle<1\).  Hence the point indexed by \(i\) is in
the relative interior of its asserted side.
\end{proof}

\begin{proposition}[Physical realization of a strict triangle cycle]
\label{prop:product-triangle-strict-physical-realization}
Suppose the active word
\[
 \sigma=(a_0,b_0,a_1,b_1,a_2,b_2)
\]
satisfies the six strict directed signs in
\eqref{eq:product-triangle-directed-signs}.  Then there is a proper
strong three-bounce \((P,Q)\)-billiard
\((q_r;p_r)_{r\in\mathbb Z/3}\) with
\begin{align}
 q_r&\in\operatorname{relint}F_P(a_{-r}),
 &p_r&\in\operatorname{relint}F_Q(b_{-r-1}),
 \label{eq:product-triangle-contact-order}\\
 q_{r+1}-q_r&\in\mathbb R_{>0}b_{-r-1}=N_Q(p_r),
 &p_{r+1}-p_r&\in-\mathbb R_{>0}a_{-r-1}=-N_P(q_{r+1}).
 \label{eq:product-triangle-strong-inclusions}
\end{align}
Its \(Q\)-length is
\begin{equation}
\label{eq:product-triangle-physical-length}
 \ell_Q(q)=\sum_r h_Q(q_{r+1}-q_r)
 =\frac1S=A_{3,\sigma}(P,Q)<A_2(P,Q).
\end{equation}
Thus the reversal from the repository active word to billiard traversal is
explicit: the billiard visits \(a_0,a_2,a_1\), and its outgoing \(Q\)-side
normals are respectively \(b_2,b_1,b_0\).
\end{proposition}

\begin{proof}
Construct a closed triangle \(T\), up to translation, with vertices \(s_i\)
labelled by the \(P\)-side contacts and edges
\begin{equation}
\label{eq:product-triangle-position-template}
 s_i-s_{i+1}=B_i
 \qquad(i\in\mathbb Z/3).
\end{equation}
Closure is \(\sum_iB_i=0\).  At \(s_i\), the incoming edge is \(B_i\),
and the outgoing edge is \(B_{i-1}\).  The strict cycle signs give
\[
 \langle a_i,B_i\rangle>0,
 \qquad \langle a_i,B_{i-1}\rangle<0.
\]
Therefore \(a_i\) is strictly maximized on \(T\) at \(s_i\).  Apply
Lemma~\ref{lem:product-triangle-homothetic-placement} to obtain
\(q_i'=z+t s_i\in\operatorname{relint}F_P(a_i)\).

Choose the harmless translation \(s_0=0\).  Then
\(s_2=B_2\) and \(s_1=-B_0\), so the reciprocal scale denominator is
\begin{equation}
\label{eq:product-triangle-position-denominator}
 \sum_i\langle A_i,s_i\rangle
 =-c_{10}+c_{22}=n_0+p_2=S.
\end{equation}
Thus \(t=1/S\).  Set \(q_r=q'_{-r}\).  From
\eqref{eq:product-triangle-position-template},
\[
 q_{r+1}-q_r=tB_{-r-1}
 =\frac{\gamma_{-r-1}}S b_{-r-1}.
\]

Construct independently a closed triangle \(R\), with a vertex \(v_j\)
labelled by each \(Q\)-side contact, by
\begin{equation}
\label{eq:product-triangle-momentum-template}
 v_{j-1}-v_j=-A_j.
\end{equation}
At \(v_j\), the incoming edge is \(-A_{j+1}\) and the outgoing edge is
\(-A_j\).  Hence
\[
 \langle b_j,-A_{j+1}\rangle>0,
 \qquad \langle b_j,-A_j\rangle<0,
\]
so \(b_j\) is strictly maximized there.  The homothetic placement lemma gives
\(p'_j=w+u v_j\in\operatorname{relint}F_Q(b_j)\).  Taking \(v_0=0\)
gives \(v_2=-A_0\), \(v_1=A_1\), and therefore
\begin{equation}
\label{eq:product-triangle-momentum-denominator}
 \sum_j\langle B_j,v_j\rangle
 =c_{11}-c_{02}=p_1+n_2=S.
\end{equation}
Thus \(u=1/S=t\).  Set \(p_r=p'_{-r-1}\).  Equation
\eqref{eq:product-triangle-momentum-template} yields
\[
 p_{r+1}-p_r=-\frac{\alpha_{-r-1}}S a_{-r-1}.
\]

Relative-interior side contacts have one-dimensional outer normal cones, so
the two last displays prove both inclusions in
\eqref{eq:product-triangle-strong-inclusions}.  All coefficients are positive,
and the template triangles are nondegenerate, so the billiard is proper.
Finally, \(h_Q(b_j)=1\) for every normalized facet vector.  Hence
\[
 \ell_Q(q)=\sum_r h_Q\left(\frac{\gamma_{-r-1}}S b_{-r-1}\right)
 =\frac1S\sum_j\gamma_j=\frac1S.
\]
Lemma~\ref{lem:product-triangle-sign-cell-algebra} identifies this with the
fixed-word QP action, and
Theorem~\ref{thm:product-triangle-wordwise-comparison} makes the comparison
strict.  The index definitions show directly that this billiard traversal is
the reverse of the displayed active Reeb word.
\end{proof}
\end{unverified}

The proposition proves global placement only for a strict Hamiltonian cycle
of two triangles.  It does not assert that the weak incident-zero construction
has relative-interior contacts or that every degenerate local cycle produces a
strong billiard.  Conversely, any proper strong three-bounce billiard whose
six contacts lie in relative interiors of triangle sides induces, after the
same reversal, a Hamiltonian word: closure by positive side-normal multiples
forces all three side normals of each triangle to occur.  Convexity then makes
the outgoing side functional decrease strictly and the incoming side
functional increase strictly, so all six transition signs are strict.  Thus
absence of a strict cycle excludes this proper side-contact type, but the
present packet makes no completeness claim about vertex-contact degeneracies.

\subsection{Exact finite evidence and its boundary}
\label{subsec:product-triangle-evidence}

The exact packet at
\begin{center}
\path{experiments/sys-datascience/methods/product-triangle-bounce-classification/}
\end{center}
is a finite falsifier and orientation check only.  On the 1,024 retained
triangle--triangle rows it found 814 \(q_0\)-fixed feasible word identities;
718 rows had at least one such word, and all 718 had every computed feasible
all-singleton action strictly below exact \(A_2\).  All 1,024 independently
compared stored two-bounce values matched the exact formula.

Its preregistered rational stress screen used seed \(20260714\) on 20,000
pairs: 19,968 were in the strict-sign stratum, 32 had a zero pairing, and
11,101 feasible all-singleton words were found.  It recorded zero strict
counterexamples and 27 exact boundary equality witnesses.  These counts
support the orientation and boundary split in the proof, but computation is
not a theorem.  In particular, the packet does not prove universal word
coverage, physical trajectory completeness, or any claim about degenerate
vertex-contact billiards.

The exact mathematical proof above depends on the triangle slack-coordinate
identity, row and column sums, the project QP convention, and the explicit
homothetic placement.  It does not use the finite packet as a proof premise.
